% Náhled k~sekci o~Kleeneho větě: regulární výraz, NFA a~DFA popisují tutéž třídu jazyků.
\begin{tikzpicture}[x=1cm, y=1cm]

    \tikzset{
        krabice/.style={draw, thick, rounded corners=3pt, minimum width=26mm,
                        minimum height=10mm, align=center, fill=maincolor!10}
    }

    \node[krabice] (re)  at (0,0)     {regulární\\ výraz};
    \node[krabice] (nfa) at (5,0)     {NFA};
    \node[krabice] (dfa) at (2.5,-2.2) {DFA};

    \draw[<->, >=stealth, very thick, darkred] (re) -- (nfa);
    \draw[<->, >=stealth, very thick, darkred] (nfa) -- (dfa);
    \draw[<->, >=stealth, very thick, darkred] (re) -- (dfa);

    \node[font=\small] at (2.5,-3.4) {tatáž třída \markdarkred{regulárních jazyků}};

\end{tikzpicture}
